%------------------------------------------------------------------------------%
\chapter*{Prefácio}

%------------------------------------------------------------------------------%
\section*{Sobre a Disciplina}

Geometria Analítica e Álgebra Linear

Carga horária: 60 h.a.

Ementa: Matrizes, sistemas de equações lineares e determinantes. Álgebra
vetorial. Retas e planos. Espaços vetoriais em R2 e R3. Autovalores e
autovetores de matrizes. Diagonalização de matrizes. Cônicas.

Bibliografia Básica
1. CAMARGO, I.; BOULOS, P. Geometria Analítica - Um Tratamento Vetorial. 3. ed. São Paulo: Prentice-Hall, 2005.
2. BOLDRINI, J. L.; et al. Álgebra Linear. 3. ed. São Paulo: HARBRA, 1986.
3. STEINBRUCH, A.; WINTERLE, P. Geometria Analítica. 2. ed. São Paulo: Makron Books,1987.

Bibliografia Complementar
1. WINTERLE, P. Vetores e geometria analítica. 2. ed. São Paulo: Pearson Education do Brasil, 2000.
2. SANTOS, R. J. Matrizes, vetores e geometria analítica. Belo Horizonte: Imprensa Universitária UFMG, 2007. http://www.mat.ufmg.br/~regi/gaalt/gaalt1.pdf
3. SANTOS, R. J. Um curso de geometria analítica e álgebra linear. Belo Horizonte: Imprensa Universitária da UFMG, 2010. http://www.mat.ufmg.br/~regi/gaalt/gaalt0.pdf
4. SANTOS, N. M., Vetores e matrizes: uma introdução à álgebra linear. 4ed. São Paulo: Thomson Learning, 2005.
5. THOMAS, George B. Cálculo. 11. ed. São Paulo: Pearson Education do Brasil Ltda., 2008. v. 2.

%------------------------------------------------------------------------------%
